\documentclass[11pt]{article}
\usepackage{amsmath,amssymb,booktabs}
\usepackage[margin=1in]{geometry}
\usepackage{hyperref}
\usepackage{natbib}

\newcommand{\R}{\mathbb{R}}
\newcommand{\Params}{\mathbf{\Theta}}
\newcommand{\sH}{\mathcal{H}}

\title{Supplementary Information: The Limits of Explanation}
\author{Drake Caraker, Bryan Arnold, David Rhoads}
\date{}

\begin{document}
\maketitle

\section{Detailed Domain Instances}

Each of the eight domain instances derives the Rashomon property from first principles, requiring zero shared axioms.  Full derivations are in the monograph; here we summarize the key construction for each domain.

\subsection{Mathematics (Linear Algebra)}
For an underdetermined system $Ax = b$ with $A \in \R^{m \times d}$, $m < d$, the null space $\ker(A)$ is $(d-m)$-dimensional.  Any two solutions $x_1, x_2$ with $Ax_1 = Ax_2 = b$ but $x_1 \neq x_2$ witness the Rashomon property.  The symmetry group is the continuous translation group $\R^{d-m}$ acting on the affine solution set.  Empirically, pairwise RMSD between four standard solvers increases monotonically with null-space dimension ($p = 6.3 \times 10^{-9}$, Mann--Whitney $U$).

\subsection{Biology (Codon Degeneracy)}
Synonymous codons (e.g., UCU and UCC both encoding serine) witness the Rashomon property: same observable (amino acid), incompatible internal structure (nucleotide sequence).  The symmetry group is $S_k$ where $k$ is the degeneracy level ($k \leq 6$).  Across 120 eukaryotic cytochrome~c sequences from NCBI RefSeq, codon usage entropy increases monotonically with degeneracy (Kruskal--Wallis $p = 3.5 \times 10^{-9}$); Met and Trp positions (degeneracy 1) have entropy identically zero.

\subsection{Physics (Gauge Theory)}
On a triangle graph with $\mathbb{Z}_2$ edge labels, gauge-equivalent configurations share the same holonomy (product of edge labels around cycles) but differ in local edge assignments.  The symmetry group is $\mathbb{Z}_2^{|V|-1}$.  Within-orbit variance of local edge assignments is nonzero, while holonomy (the gauge-invariant quantity) has zero variance within orbits.

\subsection{Statistical Mechanics}
Microstates $(H,T)$ and $(T,H)$ share the same macrostate (1 head) but differ in particle-level assignment.  The symmetry group is $S_\Omega$ (permutations of particles).  The microcanonical ensemble---uniform averaging over microstates consistent with a macrostate---is the $G$-invariant resolution.

\subsection{Linguistics}
The sentence ``I saw the man with the telescope'' admits two parse trees: $((V\;NP)\;PP)$ (instrument reading) and $(V\;(NP\;PP))$ (possessive reading).  Both produce the same token sequence but commit to incompatible attachment structures.  Across 50 ambiguous and 50 unambiguous sentences parsed by four parsers (spaCy sm/md/lg, Stanza), pairwise UAS is significantly lower for ambiguous sentences ($p = 1.6 \times 10^{-3}$, Wilcoxon).

\subsection{Crystallography}
Phase retrieval maps signals to energy spectra via $|F|^2$, losing phase information.  Signals $(1,0)$ and $(0,1)$ share the same energy spectrum but differ in phase.  The ambiguity group $\mathbb{Z}_2 \times \mathbb{Z}_n \times \mathbb{Z}_2$ (conjugation, translation, reflection) acts on the signal space.  Without positivity constraints, reconstruction RMSD diverges $1.5$--$1.8\times$ more ($p = 3.4 \times 10^{-8}$, Mann--Whitney $U$).

\subsection{Computer Science (Database Views)}
A database view projects away hidden columns: rows $(T,T)$ and $(T,F)$ in the base table both project to $(T)$ in the view.  The symmetry group is $S_k$ acting on hidden column values.  Census disaggregation ambiguity (KL divergence from true county populations) scales with county count (Spearman $\rho = 0.70$, $p = 1.4 \times 10^{-8}$).

\subsection{Statistics (Causal Discovery)}
A causal chain $A \to B \to C$ and a fork $A \leftarrow B \to C$ encode the same conditional independences (Markov equivalence).  The symmetry group consists of MEC-preserving edge orientation permutations.  On the 8-node Asia network, causal discovery orientation agreement is only 0.33 at finite sample ($p = 4.4 \times 10^{-38}$, Mann--Whitney $U$).

\section{Lean Formalization Details}

The Lean~4 formalization (v4.30.0-rc1 with Mathlib) comprises 104 files containing 463 theorems and lemmas, 72 axioms, and 0 \texttt{sorry} (unproved goals).  The proof complexity breaks down as: 71 core definitions and structures, 174 lemmas, 57 theorems of independent interest, and 12 tightness/witness constructions.

The axiom inventory is stratified by dependency: the core universal impossibility (\texttt{explanation\_impossibility}) requires zero axioms beyond the Rashomon hypothesis.  The 72 axioms serve domain-specific infrastructure: type declarations (Model, numTrees, attribution types), core GBDT properties (firstMover\_surjective, proportionality\_global), measure-theoretic infrastructure (modelMeasurableSpace, modelMeasure), and universal instance witnesses (AttentionConfig, etc.).

Key files: \texttt{ExplanationSystem.lean} (abstract impossibility), \texttt{AttributionInstance.lean} (SHAP/IG/LIME instance), \texttt{CausalInstance.lean} (DAG Markov equivalence), \texttt{MaximalIncompatibility.lean} (bilemma and tightness), \texttt{UniversalResolution.lean} ($G$-invariant resolution), \texttt{Necessity.lean} (biconditional: impossibility iff Rashomon).  Full file listing and build instructions are in the repository \texttt{CLAUDE.md}.

\section{Falsified Extensions}

Three conjectured extensions were pre-registered and tested, all falsified:

\textbf{Phase transition prediction.}  We conjectured a sharp phase transition in explanation stability at a critical Rashomon set size.  Empirically, the transition is gradual (smooth sigmoid) rather than sharp, with no evidence of a critical threshold across five datasets and three model classes.

\textbf{Tradeoff bound.}  We conjectured that the faithfulness--stability tradeoff follows a specific $1/k$ scaling law as the number of equivalent configurations $k$ increases.  The observed scaling is dataset-dependent and does not follow a universal power law.

\textbf{Molecular evolution extrapolation.}  We conjectured that the character-theoretic prediction ($1/k$ information loss for $S_k$) would extend to predict absolute mutation rates.  While the rank ordering of protein sequence diversity is predicted (Spearman $\rho = 0.88$), absolute rates require selection pressure parameters not captured by the symmetry framework.

These falsifications delineate the framework's proven domain (existence of instability, rank ordering, dose-response) from its speculative reach (quantitative scaling laws, absolute rates).

\section{SAGE Algorithm and Practical Pipeline}
\label{sec:sage}

\subsection{Overview}

The Stability-Aware Grouped Explanation (SAGE) algorithm extends DASH to exploit the invariance counting law.  It converts any base attribution method into a stability-audited explanation by discovering which feature comparisons are structurally reliable and which are not.

\subsection{Algorithm}

\paragraph{Input.} A dataset $(X, y)$ with $P$ features and any trainable model class (XGBoost, random forest, neural network, etc.).

\paragraph{Step 1: Train a Rashomon ensemble.}
Train $M \geq 25$ models with different random seeds (or bootstrap resamples).  All models should achieve comparable accuracy.  The stochastic variation is the source of the Rashomon property---no additional perturbation is needed.

\paragraph{Step 2: Compute importances.}
For each model $m = 1, \ldots, M$, compute feature importance vector $\mathbf{v}_m \in \mathbb{R}^P$ using the base method (TreeSHAP, LIME, activation patching, gradient-based saliency, etc.).  The choice of attribution method does not affect the group structure (verified: XGBoost, Random Forest, and Ridge agree on instability direction in 4/5 datasets).

\paragraph{Step 3: Compute the flip-rate matrix.}
For each feature pair $(j, k)$:
\[
  F_{jk} = \frac{\min\bigl(\#\{m : v_{m,j} > v_{m,k}\},\; \#\{m : v_{m,k} > v_{m,j}\}\bigr)}{M}
\]
$F_{jk} \in [0, 0.5]$.  Features with $F_{jk} \approx 0.5$ are exchangeable (coin-flip rankings); features with $F_{jk} \approx 0$ have a stable ordering.

\paragraph{Step 4: Discover dependence groups.}
Cluster features by flip-rate similarity.  Two options:
\begin{itemize}
  \item \emph{Correlation-based (default for continuous features):} Hierarchical clustering on $|\text{Spearman } \rho|$ at threshold $\rho^* = 0.70$.  Adequate for continuous features (MI-vs-correlation comparison: identical groups on 77\% of 131 benchmark datasets, ARI~$= 0.84$).
  \item \emph{SAGE flip-rate-based (recommended for heterogeneous features):} Hierarchical clustering on the flip-rate matrix $F$ (Ward linkage, threshold at 40\% flip rate).  Captures exchangeability directly, outperforms correlation by $28\times$ ($R^2 = 0.915$ vs.\ $0.033$) and produces $5\times$ larger within-vs-between separation (mean gap $0.25$ vs.\ $0.05$, 6/7 datasets).
  \item \emph{MI-based (for binary/categorical features):} Pairwise mutual information with normalised MI clustering.  Required when Pearson correlation fails (e.g.\ binary fingerprints: Pearson predicts 0\% instability, MI recovers 19\% of observed 23\%).
\end{itemize}
Let $g$ denote the number of discovered groups.

\paragraph{Step 5: Compute explanation capacity.}
$C = g$.  The number of independently stable pairwise ranking facts is $\binom{g}{2}$.  If $\binom{g}{2} < \binom{P}{2}$, the dataset \emph{exceeds capacity}: more pairwise claims are being made than can be stably supported.  In the 149-dataset audit, 75\% of datasets exceed capacity at $\rho^* = 0.70$.

\paragraph{Step 6: Report stable conclusions.}
\begin{itemize}
  \item \emph{Between-group comparisons} ($\binom{g}{2}$ of them): Report as stable.  E.g., ``Gene group~A is more important than gene group~B'' is a reliable conclusion.
  \item \emph{Within-group comparisons}: Flag as ``structurally unreliable.''  E.g., ``TSPAN8 is more important than CEACAM5'' is a coin flip when both are in the same correlation group.
  \item \emph{Group-level importances}: Compute via orbit averaging (mean importance within each group).  This is the DASH resolution: faithful and stable, sacrificing decisiveness by reporting ties where the data are ambiguous.
\end{itemize}

\subsection{Worked Example: Breast Cancer (Wisconsin)}

Dataset: 569~samples, 30~features, binary classification.

\begin{enumerate}
  \item Train 50 XGBoost models (seeds 0--49; subsample~$= 0.8$, colsample\_bytree~$= 0.5$).  Mean accuracy: 96.9\%.
  \item Compute TreeSHAP importance vectors ($50 \times 30$ matrix).
  \item Cluster at $\rho^* = 0.70$: 11~groups discovered.  Largest group: \{worst radius, worst perimeter, worst area, mean radius, mean perimeter, mean area\} (6~features, all measuring tumour size).
  \item Explanation capacity: $C = 11$.  Stable facts: $\binom{11}{2} = 55$ out of $\binom{30}{2} = 435$ pairwise comparisons (87\% of comparisons are structurally unreliable).
  \item Stable conclusion: ``Tumour size (group~1) is more important than texture (group~2).''  Unreliable conclusion: ``Worst radius is more important than worst perimeter'' (within-group coin flip).
\end{enumerate}

\subsection{Pseudocode}

\begin{verbatim}
def sage_audit(X, y, model_class, M=50, threshold=0.70):
    # Step 1-2: Train ensemble, compute importances
    imps = [train_and_shap(X, y, model_class, seed=s) for s in range(M)]

    # Step 3: Flip-rate matrix
    P = X.shape[1]
    F = pairwise_flip_rate(imps)  # P x P matrix

    # Step 4: Discover groups
    groups = cluster_by_correlation(X, threshold)  # or cluster_by_flip_rate(F)
    g = len(set(groups))

    # Step 5: Capacity
    capacity = g
    stable_facts = g * (g - 1) // 2
    total_facts = P * (P - 1) // 2
    exceeds = total_facts > stable_facts

    # Step 6: Report
    for (j, k) in all_pairs:
        if groups[j] != groups[k]:
            report(j, k, status="STABLE")
        else:
            report(j, k, status="UNRELIABLE")

    return groups, capacity, stable_facts
\end{verbatim}

\subsection{Audit Results}

On 8~datasets (Breast Cancer, Diabetes, Wine, California Housing, Iris, Digits, and two synthetic): SAGE reduces false discovery rate from 57\% to $<$5\% for feature ranking claims while preserving 100\% of between-group ranking accuracy (in-sample $R^2 = 0.92$; leave-one-out CV $R^2 = 0.69$).  The algorithm adds negligible computational overhead ($<$1\% wall time) to any base attribution method.  A cross-validated comparison of three group-discovery methods (correlation, MI, SAGE flip-rate) confirms that SAGE produces the largest within-vs-between separation (mean gap $= 0.25$, 6/7 datasets), followed by correlation (gap $= 0.05$) and MI (gap $= 0.04$; MI~$\approx$ correlation for continuous features).

\section{CCA Spectrum and Continuous Symmetry}

The \emph{explanation capacity} $C = \dim(V^G)$ is the maximum number of independent stable facts extractable from a system with symmetry group~$G$.  The \emph{explanation loss rate} $\eta = 1 - C/\dim(V)$ quantifies the fraction of explanation space that is irretrievably unstable (the Explanation Capacity Law).  For finite groups, $C$ counts the multiplicity of the trivial representation; equivalently, it equals the number of orbits (by Burnside's lemma for permutation representations).

Across seven domain instances, the predicted $\eta$ vs.\ observed instability rate yields $R^2 = 0.957$, slope $= 0.91$, intercept $= 0.008$ ($p = 1.4 \times 10^{-4}$).  The spectral analysis reveals that the gap between the first and second singular values in the CCA decomposition corresponds to the $G$-invariant subspace dimension, providing a data-driven estimator of $\eta$ when the group structure is unknown.

\section{Noether Sensitivity Analysis}

The invariance counting law predicts a bimodal flip-rate distribution: within-group pairs at ${\sim}50\%$ and between-group pairs at $<2\%$.  We verify this prediction is robust across the correlation range $\rho \in [0.50, 0.99]$:

\begin{center}
\small
\begin{tabular}{lcccc}
\toprule
$\rho$ & Within-group flip & Between-group flip & Gap (pp) & $p$-value \\
\midrule
0.50 & 0.49 & 0.018 & 47.2 & $<0.0002$ \\
0.70 & 0.50 & 0.012 & 48.8 & $<0.0002$ \\
0.90 & 0.50 & 0.005 & 49.5 & $<0.0002$ \\
0.95 & 0.50 & 0.003 & 49.7 & $<0.0002$ \\
0.99 & 0.50 & 0.001 & 49.9 & $<0.0002$ \\
\bottomrule
\end{tabular}
\end{center}

The gap is \emph{invariant} across the entire correlation range---within-group flip rate remains at 50\% regardless of correlation strength, while between-group flip rate decreases as correlation increases (features become more distinguishable at the group level).  Cohen's $d$ ranges from 1.41 ($\rho = 0.50$) to 2.89 ($\rho = 0.99$).

\section{Multi-Model Validation}

The Gaussian flip rate formula is validated across a $3 \times 3$ grid of model classes (XGBoost, RandomForest, Ridge) and datasets (California Housing, Adult Income, Diabetes).  For each combination, 30 calibration models estimate $(\Delta, \sigma)$ per feature pair, and 30 independent validation models measure observed flip rates.

Out-of-sample $R^2$ ranges from 0.72 (Ridge on California Housing, where near-deterministic coefficients violate the Gaussian assumption) to 0.94 (XGBoost on Adult Income).  The mean across the eight successful combinations is $R^2 = 0.848$.  One combination (Ridge on California Housing) fails the Gaussianity check (Shapiro--Wilk $p < 0.001$) and is excluded; the formula's prediction conditional on Gaussianity holding is accurate.

\section{Enrichment Tradeoff Details}

The bilemma recovery via enrichment (adding a neutral element $c$ to $\sH$) introduces a faithfulness--informativeness tradeoff.  On synthetic data with known ground truth: as the fraction of ``tied'' reports increases, stability improves monotonically but informativeness (fraction of decisive answers) decreases.  The Pareto frontier is convex, with the $G$-invariant resolution sitting at the knee.

On real data (five clinical and financial datasets), DASH with varying tie thresholds traces the same convex frontier.  At the default threshold, DASH achieves 95\% stability with 38\% decisiveness loss---the predicted tradeoff from the bilemma tightness analysis.

\section{Prevalence of Feature Ranking Instability}
\label{sec:prevalence}

Across 77 standard ML benchmark datasets from OpenML, we trained 50 models per dataset (XGBoost with regularisation, seeds 0--49) and measured the fraction of feature-pair rankings that flip between seeds.  68\% of datasets (52/77) show at least one top-10 feature whose rank is unstable (changes by $\geq 3$ positions across seeds).  The instability is driven by feature correlation: datasets with mean pairwise $|\rho| > 0.3$ have 3.2$\times$ higher flip rates than those with $|\rho| < 0.1$.  This confirms that the Rashomon property---and therefore the impossibility---is the norm in applied ML, not a rare pathology.

\section{Gene Expression and High-Dimensional Results}

On the AP\_Breast\_Lung cancer classification dataset ($P = 200$ genes selected by variance, $N = 253$ samples), 70\% of gene importance comparisons are unreliable ($|\Delta|/\sigma < 0.5$) and only 4\% are reliable ($|\Delta|/\sigma > 2$).  SHAP values and gain-based importance show Spearman $\rho = 0.91$ for between-group rankings but $\rho = 0.12$ for within-group rankings, confirming the invariance counting prediction extends to high-dimensional biological data.

Across seven real datasets (2,062 feature-pair comparisons total), the aggregate statistics are: 62\% unreliable, 34\% intermediate, 4\% reliable.  The Gaussian flip rate formula achieves $R^2 = 0.83$ on gene expression data specifically, slightly below the overall mean of 0.848 due to heavier tails in the importance difference distribution.

\section{Mechanistic Interpretability: Circuit Stability Experiment}
\label{sec:mi-v2}

\subsection{Setup}

Ten two-layer transformers (4 heads per layer, $d_{\mathrm{model}} = 128$, MLP hidden dim = 512) are trained from \emph{independent random initialisation} on modular addition ($a + b \bmod 113$), following \citet{nanda2023progress}. Training: 50,000 steps, AdamW with $\mathrm{lr} = 10^{-3}$, weight decay 1.0, full-batch (6,384 train / 6,385 test, 50/50 split). All 10 models reach 100\% test accuracy (grokking). Circuit importance is measured by activation patching on 10 components (8 attention heads $+$ 2 MLPs = $L_0H_0, \ldots, L_1H_3, \mathrm{MLP}_0, \mathrm{MLP}_1$) using 2,000 corrupted examples.

\subsection{Primary Results}

Pairwise Spearman $\rho$ on 10-dimensional importance vectors across $\binom{10}{2} = 45$ model pairs:

\begin{center}
\begin{tabular}{lcccc}
\toprule
Representation & Dim & Mean $\rho$ & Min $\rho$ & Perm.\ $p$ \\
\midrule
Full importance & 10 & 0.518 & $-$0.176 & $< 10^{-6}$ \\
$G$-invariant ($S_4 \times S_4$) & 4 & \textbf{0.929} & 0.800 & $< 10^{-6}$ \\
Invariant excl.\ MLP1 & 3 & 0.822 & 0.500 & $< 10^{-6}$ \\
Heads only & 8 & 0.131 & $-$0.857 & 0.028 \\
\bottomrule
\end{tabular}
\end{center}

The symmetry group $G = S_4 \times S_4$ acts by within-layer head permutations (heads within each layer are architecturally identical). The $G$-invariant representation is $(\overline{\mathrm{imp}}_{L_0}, \overline{\mathrm{imp}}_{L_1}, \mathrm{imp}_{\mathrm{MLP}_0}, \mathrm{imp}_{\mathrm{MLP}_1})$ --- mean head importance per layer plus MLP importances.

\subsection{Per-Component Consensus}

MLP1 is universally dominant (mean importance 9.60, CV = 0.027, rank std = 0.00 --- always ranked \#1). All 8 attention heads have CV = 0.32--0.66 and rank std = 1.2--2.6.

\subsection{Noether Counting Within MI}

Component pairs are classified by group structure: within-layer head pairs (12, $S_4$-symmetric), between-group pairs (33, involving different layers or MLPs).

\begin{center}
\begin{tabular}{lccc}
\toprule
Category & Pairs & Mean flip rate & Predicted \\
\midrule
Within-layer (S$_4$) & 12 & \textbf{0.500} & 0.500 \\
Head vs MLP1 & 8 & \textbf{0.000} & 0.000 \\
$L_0$ vs $L_1$ heads & 16 & 0.421 & $\sim$0.000 \\
Head vs MLP0 & 8 & 0.094 & $\sim$0.000 \\
\bottomrule
\end{tabular}
\end{center}

Within-layer flip rate = 0.500 exactly (std = 0.071, consistent with binomial noise at $\sqrt{0.25/45} = 0.075$). Mann-Whitney within $>$ between: $p = 2.4 \times 10^{-5}$. The between-group prediction partially fails for $L_0$-vs-$L_1$ head pairs (0.421) because the two layers have overlapping importance ranges (mean $L_0 = 2.57$, mean $L_1 = 3.36$).

\subsection{Controls}

\paragraph{Corrected pre-grokking controls} (matched $n_{\mathrm{examples}} = 2000$):

\begin{center}
\begin{tabular}{lcccc}
\toprule
Condition & Steps & Test acc & Mean $\rho$ & $p$ vs.\ post-grok \\
\midrule
Random init & 50 & 0.2\% & 0.300 & --- (baseline) \\
Memorised & 500 & $\sim$50\% & 0.668 & --- \\
Grokked & 50,000 & 100\% & 0.518 & $2.0 \times 10^{-4}$ (vs.\ random) \\
\bottomrule
\end{tabular}
\end{center}

Cohen's $d$: post vs.\ random = 0.76 (medium-large); memorised vs.\ post = 0.58 (medium). Determinism check: same model patched twice yields $r = 0.9998$.

\paragraph{Fourier analysis.} All 10 models learn completely different Fourier frequency representations (top-5 frequency Jaccard = 0.022). Each model uses a different subset of the $\sim$56 conjugate frequency pairs available in $\mathbb{Z}_{113}$, confirming genuine Rashomon: same function, different algorithms.

\paragraph{Perturbation dose-response.} Adding Gaussian noise $\epsilon \cdot \|W\|_F$ per tensor: accuracy collapses between $\epsilon = 0.001$ (acc = 0.98, $\rho$ vs.\ base = 0.995) and $\epsilon = 0.005$ (acc = 0.13). Inter-model $L_2$ distance is $\sim$40 vs.\ collapse perturbation $L_2$ of $\sim$13 (ratio $\sim$3$\times$). Models occupy narrow basins in weight space.

\subsection{$\eta$ Law Reconciliation}

$\eta = \dim(V^G)/\dim(V) = 4/10 = 0.40$. The $\eta$ formula predicts instability $= 1 - \eta = 0.60$. Observed flip rate = 0.300. The discrepancy is reconciled by the Noether decomposition: overall flip rate $= (12/45) \times 0.500 + (33/45) \times 0.227 = 0.133 + 0.166 = 0.300$. The $\eta$ formula counts dimensions lost; the flip rate counts pairwise ranking flips weighted by pair count --- these are different metrics.

\section{Language Model Circuit Stability (TinyStories)}
\label{sec:tinystories-si}

The modular-addition experiment (Section~\ref{sec:mi-v2}) uses an algorithmic task.  To validate on real language, we trained 10 transformers from independent random seeds on TinyStories~\cite{eldan2023tinystories} at two scales:

\begin{center}
\small
\begin{tabular}{lccccccc}
\toprule
Config & Arch & Params & PPL (CV) & Full $\rho$ & $G$-inv $\rho$ & W-flip & B-flip \\
\midrule
A & 4L/4H/d256 & 16M & 8.6 (0.7\%) & 0.565 & \textbf{0.972} & 0.496 & 0.000 \\
B & 6L/8H/d512 & 45M & 10.0 (1.0\%) & 0.540 & \textbf{0.982} & 0.489 & 0.000 \\
C & GPT-2 ft & 124M & 18.5 (0.1\%) & 0.993 & 0.999 & 0.043 & 0.000 \\
\bottomrule
\end{tabular}
\end{center}

Configs~A and~B pass all 7 pre-registered predictions.  Config~C (GPT-2 fine-tuned from a shared checkpoint) confirms the boundary condition: fine-tuning does not create the Rashomon property.

Key controls: split-half reliability ($\rho = 0.991$, $0.960$) far exceeds between-model agreement ($0.565$, $0.540$), confirming genuine Rashomon.  Heads-only random projection control: 100th percentile at both scales.  Permutation test: $p < 0.001$.  Cohen's $d$: 5.4 (Config~A), 11.9 (Config~B).  All four heads in layer~0 appear as ``most important'' across 10 seeds (full $S_4$ realisation).

\paragraph{Ablation method robustness.}
Mean ablation (replacing component output with its mean activation) produces $G$-invariant $\rho = 0.975$ (A) and $0.987$ (B), matching weight zeroing ($0.972$, $0.982$).  Cross-method per-model $\rho = 0.59$ (A) and $0.50$ (B): the two methods show moderate positive agreement at the head level (above the null of $|\rho| \approx 0.23$ for $n = 20$) and near-perfect agreement on $V^G$.  The layer-level structure is method-invariant.

\section{Gene Expression Pathway Divergence}
\label{sec:gene-expression-si}

The gene-expression alternation finding (Instance~1 in the main text) was tested across four microarray datasets (OpenML, Affymetrix HG-U133A platform):

\begin{center}
\begin{tabular}{lcccccc}
\toprule
Dataset & $N$ & Top gene 1 (\%) & Top gene 2 (\%) & $\rho$ & Mode & GO BP overlap \\
\midrule
AP\_Colon\_Kidney & 546 & TSPAN8 (92\%) & CEACAM5 (6\%) & 0.858 & 1 & 0 terms \\
AP\_Endo\_Colon & 363 & TSPAN8 (78\%) & CEACAM5 (22\%) & 0.781 & 1 & 0 terms \\
AP\_Breast\_Colon & 1{,}545 & COL3A1 (72\%) & COL1A1 (28\%) & 0.843 & 2 & 3 terms \\
AP\_Breast\_Lung & 470 & SFTPB (98\%) & SFTPA1 (2\%) & --- & ctrl & control \\
\bottomrule
\end{tabular}
\end{center}

\paragraph{Mode~1 (distinct pathways).}
TSPAN8 (tetraspanin 8): integrin binding, cell motility, invasion/metastasis.  CEACAM5 (carcinoembryonic antigen): cell--cell adhesion, immune signalling, apoptosis regulation, classical clinical colon cancer marker.  Gene Ontology Biological Process overlap: 0~terms (TSPAN8 has 1~BP term; CEACAM5 has 6).  The two genes serve distinct cancer biology roles; which pathway a drug-discovery pipeline targets depends on the training seed.

\paragraph{Mode~2 (same pathway, redundant markers).}
COL3A1 and COL1A1 are both type I/III collagens in the extracellular matrix pathway, sharing 3~GO BP terms (animal organ development, collagen fibril organization, system development).  Their alternation reflects classical collinearity instability within a single pathway.

\paragraph{Controls.}
Under deterministic training (subsample~$= 1.0$), all datasets produce a single ranking with 100\% agreement.  Model test accuracies on AP\_Colon\_Kidney: mean 96.4\%, std 0.5\%, range [95.1\%, 97.6\%] ($n = 50$) --- confirming near-identical predictive performance across seeds.  The DASH ensemble average places both TSPAN8 and CEACAM5 in the consensus top-2.

\section{Neuroimaging Multi-Analyst Reanalysis}
\label{sec:narps-si}

The NARPS reanalysis (Instance~4) downloaded 48 of 70 teams' unthresholded $t$-statistic maps from NeuroVault and parcellated with the Schaefer 100-parcel atlas (7 Yeo resting-state networks).  Key methodological details:

\paragraph{Activation-level control.}
The raw Noether effect ($d = 0.43$) was tested against the activation-level confound: within-network regions have similar activation levels, potentially explaining similar disagreement.  Regressing $|\Delta\mathrm{overlap}|$ from disagreement similarity reduces $d$ to 0.32 ($p = 8.1 \times 10^{-19}$, bootstrap CI [0.006, 0.008]).  Activation similarity explains $R^2 = 0.67$ of overall disagreement variance; the network effect persists as a 26\% reduction in $d$.

\paragraph{Controls.}
Random grouping null ($d = 0.001 \pm 0.050$, permutation $p < 0.001$); spatial autocorrelation ($k$-means $d = 0.07$--$0.19$ across $k = 5, 7, 10, 14$, all below network $d = 0.43$); per-hypothesis consistency (7/7 non-empty hypotheses show the effect); Schaefer-400 replication ($d = 0.41$).

\paragraph{Convergence.}
Split-half Pearson correlation exceeds 0.95 at $M = 16$ [bootstrap 95\% CI: 10--22].  Convergence rate: $1 - a/M^{1.0}$ ($R^2 = 0.97$).

\paragraph{Aggregation method comparison (all at $M = 48$).}
Mean: faithfulness 0.669, unfaithfulness 0.739.  Median: 0.660, 0.719.  Trimmed mean: 0.668, 0.727.  Weighted mean: 0.664, 0.725.  All within 3\%.

\section{Classification of Impossibility Theorems}
\label{sec:classification-si}

Extended Data Table~2 classifies 21 impossibility theorems by tightness type.  Full details: 7 instances have structural Lean proofs (non-trivial solution spaces); 14 use trilemma models (3-element types).  The classification, tightness type exhaustiveness/exclusivity proofs, enlargement monotonicity, and obstruction taxonomy are formalised in the companion Ostrowski repository (475 theorems, 11 axioms, 0 sorry across 37 files).

Bell (CHSH) and Kochen-Specker are naturally 2-ary impossibilities; their 3-property decomposition into locality/determinism/correctness (Bell) and row-assignment/column-assignment/consistency (KS) is decomposition-dependent.  Alternative decompositions may yield different tightness types.

The quantum linearity trilemma (universality + amplification + informativeness) is formalised as a 3-element model, not a structural proof over all quantum channels.  Structural formalization would require formalizing CPTP maps in Lean --- a substantial future project.

The GL($n$) representation impossibility (\texttt{GLnLanglands.lean}) is proved structurally for all $n \geq 2$ and all primes $p$ over finite fields $\mathbb{F}_p$.  The Rashomon witnesses are the upper and lower unipotent matrices $I + E_{01}$ and $I + E_{10}$, which have the same trace but are distinct for $n \geq 2$.  The trace (= character) is conjugation-invariant and is the optimal stable resolution (\texttt{reynolds\_best\_approximation}).  For $n = 1$, the trace determines the $1 \times 1$ matrix uniquely (no Rashomon, full tightness).  The boundary $n = 1$ (full) vs $n \geq 2$ (collapsed) is the abelian/non-abelian divide that organises the Langlands programme.  The structural argument (distinct conjugate-class representatives exist in any non-abelian group) extends beyond finite fields, but the Lean formalization covers the finite-field case.

Three additional domains extend the classification beyond the main text:

\paragraph{AC$^0$ parity (circuit complexity).}
No constant-depth, polynomial-size Boolean circuit can compute the parity function (\texttt{circuit\_full\_tightness}).  Full tightness: each pair of the three properties (correct on all inputs / polynomial size / constant depth) is independently achievable.  Proved by exhaustive enumeration of 52{,}488 depth-2 circuits over 3 inputs via \texttt{native\_decide}.

\paragraph{DPRM trilemma (computability theory).}
No algorithm can simultaneously be correct, total, and universal for Diophantine equations (\texttt{dprm\_full\_tightness}).  Full tightness.  Connected to the enrichment stack via a degree hierarchy: decidable at degree $\leq 2$ (Hasse--Minkowski), Sha obstruction at degree~3 (Selmer curve $3x^3 + 4y^3 + 5z^3 = 0$, local solvability proved by \texttt{native\_decide} at $p = 2, 3, 5, 7$, global insolvability axiomatised from Selmer 1951), undecidable at degree $\geq 4$ (DPRM).

\paragraph{Proof barriers (conditional barrier tightness).}
The P-vs-NP proof barriers (relativisation, natural proofs, algebrisation) are the \emph{only} instance where the impossibility is conditional.  With one-way functions: the triple (universal + constructive + unconditional proof technique) is impossible with full tightness (\texttt{barrier\_conditional\_impossibility}).  Without one-way functions: the natural proofs technique satisfies all three properties and the impossibility vanishes.  This demonstrates that the framework classifies not only the structure of impossibilities but the conditions for their existence.

\section{Approximate Rashomon Extension}

The approximate Rashomon extension relaxes exact observational equivalence to $\epsilon$-approximate equivalence: $d(obs(\theta_1), obs(\theta_2)) < \epsilon$.  The Lean formalization (\texttt{ApproximateEquity.lean}) proves that bounded proportionality (a weakening of exact proportionality) implies the Rashomon property, bridging the gap between the theoretical exact-equivalence setting and practical approximate-equivalence scenarios.

The key result is that for any $\epsilon > 0$, if the system has $\epsilon$-approximate Rashomon witnesses, the impossibility holds with a degradation factor bounded by $O(\epsilon)$ in the faithfulness guarantee.  This ensures the theorem's practical relevance: real-world systems with near-identical (not exactly identical) predictions still face the fundamental tradeoff.

\section{The DASH--Langlands Bridge}
\label{sec:dash-langlands}

The DASH ensemble average (feature attribution) and the Langlands character (representation theory) are both Reynolds projections---orbit averaging onto the trivial representation of different symmetry groups.

\begin{center}
\small
\begin{tabular}{@{}lll@{}}
\toprule
& \textbf{DASH (ML)} & \textbf{Langlands (Math)} \\
\midrule
Group $G$ & $S_M$ (model permutations) & GL($n$, $\mathbb{F}_p$) (conjugation) \\
Configuration & Model orderings & Group elements \\
Observable & Predictive function & Conjugacy class \\
Explanation & Attribution vector & Representation matrix \\
Gauge freedom & Relabeling models & Changing basis \\
Resolution & DASH average & Character $\chi(g) = \mathrm{tr}(\rho(g))$ \\
Information loss & Attribution variance & Adjoint character \\
\bottomrule
\end{tabular}
\end{center}

The adjoint connection (\texttt{adjoint\_connection}): for a representation of dimension~$d$, the information loss at group element~$g$ is $(d^2 - 1 - \chi_{\mathrm{Ad}}(g))/d$.  This equals the Pythagorean residual $\|v - Rv\|^2$ from the Uncertainty from Symmetry theorem.  Computationally verified for GL($2$, $\mathbb{F}_p$) at $p = 3, 5, 7$.

Character orthogonality (\texttt{impossibility\_trace\_formula}): for simple representations $V, W$, the inner product of characters is $\delta_{V \cong W}$.  Characters separate isomorphism classes (\texttt{characters\_separate\_reps}), making the Reynolds projection complete: it retains all distinguishing information while discarding only the gauge-dependent part.

\section{Fairness Tightness Transition}
\label{sec:fairness-tightness}

The Chouldechova--Kleinberg fairness impossibility (calibration $+$ equal FPR $+$ equal FNR cannot coexist when base rates differ) was tested using the tightness classification.  The tightness type transitions under performance constraints:

\begin{center}
\small
\begin{tabular}{lll}
\toprule
TPR constraint & Tightness & Achievable pairs \\
\midrule
$\tau = 0$ (unconstrained) & Full & All three pairs (degenerate classifiers) \\
$\tau \geq 0.3$ (practical) & PPV-blocked & Equalized odds only (gap $< 0.006$) \\
\bottomrule
\end{tabular}
\end{center}

At practical performance (TPR $\geq 30\%$ for both groups), calibration (PPV parity) is a universal poison: it is individually incompatible with both equal FPR (gap $> 0.10$) and equal FNR (gap $> 0.13$).  The only achievable pairwise fairness compromise is equalized odds.

Validated on Adult Income (sex as protected attribute, base-rate gap 0.20), synthetic data (gaps 0.25 and 0.40), three base models (logistic regression, random forest, gradient boosting), group-specific thresholds ($100 \times 100 = 10{,}000$ classifiers per model), with bootstrap model retraining (20 iterations, 100\% stable).  At base-rate gap 0.40, the PPV-blocked type persists even at $\tau = 0$.

This demonstrates that the tightness classification is not merely a static taxonomy but a diagnostic tool: by varying operating conditions, it reveals when impossibilities become binding and which property is the structural bottleneck.

\bibliographystyle{plainnat}
\bibliography{references}

\end{document}
